\documentstyle[%


righttag%


,11pt%


,amssymb%


,russian%               remove in Eng.


,izvru%                 Set izven% in Eng.


]{amsart}





%





\newcommand{\re}{\operatorname{Re}}


\newcommand{\im}{\operatorname{Im}}


\newcommand{\tts}[1]{}





\theoremstyle{definition}


\theorembodyfont{\rm}


\newtheorem{notenonum}{\hskip\parindent Замечание}


\renewcommand{\thenotenonum}{}





%\hoffset=-15mm%                  For Eng.


\setcounter{page}{38}


\god{2003}


\nomer{10~(497)} %                        Remove in Eng.


%\nomer{Vol.\,47, No.\,10}%               Set in Eng.


%\str{\pageref{begin}--\pageref{end}}%  Set in Eng.


\udk{514.752}





\author{\it А.В.\,Лобода, А.С.\,Ходарев}


% NB:   ^^ remove in Eng.


\title{Об одном семействе аффинно-однородных


вещественных гиперповерхностей


3-мерного комплексного пространства}


\thanks{Работа выполнена при поддержке программы 


Министерства образования Российской Федерации ``Университеты России


--- фундаментальные исследования'' (проект \No\,04-01-41) и Российского фонда 


фундаментальных исследований (проект 01-01-00594).}





\date{}


%\pagestyle{myheadings}%         Set in Eng.


\pagestyle{plain} %              Remove in Eng.


\begin{document}


%\thispagestyle{plain}%          Set in Eng.


\maketitle


\label{begin}








\section*{Введение}





Во многих разделах современной математики возникают задачи, связанные 


с описанием многообразий, 


однородных относительно действий различных групп 


преобразований. 


Простейшей из таких задач можно считать описание 


плоских кривых, однородных относительно аффинных преобразований плоскости.


Решение этой задачи, а также список аффинно-однородных


пространственных кривых давно известны [1]. 


При этом полный список 


аффинно-однородных поверхностей трехмерного вещественного пространства 


был получен лишь недавно в [2].





В аналогичной проблеме описания вещественных


гиперповерхностей комплексных прост\-ранств, являющихся орбитами


тех или иных групп преобразований (аффинных, голоморфных, дробно-линейных),


пока получены лишь некоторые частные решения [3]--[5]. 





В соответствии с серией работ [4]--[11] задачу описания 


однородных подмногообразий различных пространств удобно решать, используя


канонические (нормальные) уравнения таких многообразий. 


Например, уравнение 


вещественно-аналитической гиперповерхности $ M $ пространства 


$ \Bbb C^3 $, строго псевдо-выпуклой в некоторой своей точке,


можно привести комплексными линейными преобразованиями 


к виду 


\begin{equation}


v = ( | z_1 |^2 + |z_2 |^2 ) + ( \alpha_1 z_1^2 + \alpha_2 z_2^2 ) + 


\overline {( \alpha_1 z_1^2 + \alpha_2 z_2^2 ) } + 


\sum_{k \ge 3} F_k(z, \ov z, u ).


\end{equation}


Здесь $ z_1$, $z_2$, $w $ --- координаты в 3-мерном комплексном 


пространстве, $ u = \re w$, $v = \im w $; $ F_k $ --- однородный многочлен 


суммарной степени $ k $ по $ z$, $ \overline z $; 


$ \alpha_1 $ и $ \alpha_2 $ --- вещественные неотрицательные числа.





Опираясь на уравнения вида (1), мы изучаем аффинную однородность 


вещественных гиперповерхностей в пространстве $ \Bbb C^3 $.








Везде ниже будем рассматривать только жесткие поверхности, уравнения


(1) для которых свободны от вещественной переменной $ u = \re w $. 


К аффинно-однородным поверхностям такого типа относятся, например, 


трубчатые поверхности (трубки) вида 


$$ 


M = \Gamma + i\Bbb R^3


$$


с однородными относительно вещественных аффинных преобразований основаниями 


$ \Gamma \subset \Bbb R^3 $.


При этом для всех трубок, как нетрудно убедиться, коэффициенты 


$ \alpha_1$, $\alpha_2 $ равны $ 1/2 $. 











Отметим, что за счет 


несложных возмущений трубчатых многообразий 


можно построить большое семейство аффинно-однородных поверхностей, в 


нормальном уравнении (1) для которых лишь один из двух 


коэффициентов $ \alpha_1$, $\alpha_2 $ равен $ 1/2 $.





В данной статье рассматривается случай 


$ \alpha_2 = 0 < \alpha_1 \ne 1/2 $ и 


изучаются


аффинно-однородные поверхности с дискретными группами изотропии.


С учетом описания [11] групп, включающих дискретные группы,


получаем явный полный список аффинно-однородных поверхностей


с представлениями вида


\begin{equation}


 v = ( | z_1 |^2 + |z_2 |^2 ) + \alpha ( z_1^2 + \overline z_1^2 )


 + \sum_{k \ge 3} F_k(z,\overline z), \quad \alpha \ne 1/2. 


\end{equation}








\section{Алгебры векторных полей на однородных поверхностях}








Рассмотрим вещественно-аналитическую гиперповерхность 


$ M \in \Bbb C^3 $,


заданную вблизи начала координат уравнением 


$$


 \Phi(z,\overline z ) = - v + F(z, \overline z) = 


 - v + ( | z_1 |^2 + |z_2 |^2 ) + \alpha (z_1^2 + \overline z_1^2 ) +


 \sum_{k \ge 3} F_k(z, \ov z) = 0.


$$


Аффинная однородность $ M $ понимается как существование


в группе $\mathrm{Aff}\,(3, \Bbb C ) $ аффинных преобразований 3-мерного 


комплексного пространства


некоторой (локальной) подгруппы Ли $ G(M) $, транзитивно действующей на 


поверхности $ M $.








От группы Ли переходим к инфинитезимальным преобразованиям поверхности,


рассматриваемой вблизи начала координат.


В современной терминологии это означает рассмотрение алгебры векторных полей 


(касательных к многообразию), соответствующей группе $ G(M) $. 


Из однородности поверхности $ M $ 


вытекает наличие большой (как минимум 5-мерной) алгебры таких полей на $ M $.





В силу линейного характера используемых аффинных преобразований 


соответствующие векторные поля также являются линейными. Каждое из них 


можно представить в виде


\begin{multline}


Z = ( A_1 z_1 +  A_2 z_2 + A_3 w + p_1 )


{\partial}_{z_1} + ( B_1 z_1 + 


 B_2 z_2 + B_3 w + p_2 ){\partial}_{z_2} +


\\


+( C_1 z_1 +  C_2 z_2 + C_3 w + q ){\partial}_{w},


\end{multline}


где $ A_k$, $B_k$, $C_k$, $p_1$, $ p_2$, $q $ --- комплексные константы.





Имея линейное векторное поле на однородной вещественной поверхности, 


можно рассмотреть тождество


\begin{equation}


 \re \{ Z(\Phi) \}_{|M} \equiv 0,


\end{equation}


означающее факт касания этим полем поверхности $M$.





Определяя вес монома 


$ z_1^k z_2^l \overline z_1^s \overline z_2^t u^m $


как сумму $ k + l + s + t + 2m $, рассмотрим компоненты 


весов 0, 1, 2, 3 тождества (4). Соответствующие 


уравнения имеют вид


\begin{align}


\text{вес 0 }: \quad 


& \re \{ 


q \tfrac i2 \} \equiv 0 , \\


%


\text{вес 1 }: \quad 


&\re \{ 


p_1 (\overline z_1 + 2 \alpha z_1 ) +


p_2 \overline z_2 + ( C_1 z_1 + C_2 z_2) \tfrac i2 \} \equiv 0 , \\


%


\text{вес 2 }: \quad


&\re \{ 


 (A_1 z_1 + A_2 z_2 ) 


 {(F_2)}_{ z_1}


 + p_1 { (F_3)}_{ z_1} + \notag \\


 %


&\quad + 


 (B_1 z_1 + B_2 z_2 ) 


 { (F_2)}_{z_2} 


 + p_2 { (F_3)}_{z_2} + 


 C_3 (u + i F_2 ) \tfrac i2 \} \equiv 0 , \\


%


\text{вес 3 }: \quad


& \re  \{ 


 (A_1 z_1 + A_2 z_2 ) 


 {(F_3)}_{z_1} +


 A_3 (u + i F_2 ) {(F_2)}_{ z_1} +


 p_1 {(F_4)}_{ z_1} + \notag \\


%


&\quad + 


 (B_1 z_1 + B_2 z_2 ) 


 {(F_3)}_{ z_2} +


 B_3 (u + i F_2 ) {(F_2)}_{ z_2} +


 p_2 {(F_4)}_{z_2} + \notag \\


%


&\quad +  ( - C_3/2 ) F_3 \} \equiv 0.


\end{align}


Из самого простого тождества (5) сразу следует, что


$ q \in \Bbb R $, а из (6) выводятся два равенства: 


$$


 C_1 = 2 i ( \overline p_1 + 2 \alpha_1 p_1), \quad


 C_2 = 2 i \overline p_2. 


$$





Для получения более точной информации об однородных поверхностях 


естественно использовать алгебраическую структуру множества


линейных векторных полей вида (3), касательных к изучаемой поверхности.


Запишем всякое такое поле в виде квадратной матрицы 4-го порядка


$$


Z=


\begin{pmatrix}


A_1 & A_2 & A_3 & p_1 \\


B_1 & B_2 & B_3 & p_2 \\


C_1 & C_2 & C_3 & q \\


0 & 0 & 0 & 0


\end{pmatrix}


,


$$


обозначив ее той же буквой $ Z $, что и само поле.


Тогда множество матриц, соответствующих линейным полям на 


аффинно-однородной поверхности $ M $, образует 


алгебру Ли относительно обычной матричной скобки 


$$


 [ Z_1, Z_2 ] = Z_1 Z_2 - Z_2 Z_1. 


$$





Исходя из этих соображений с учетом изучения компонент второго и третьего


веса основного


тождества (4), приходим к следующему результату.








\begin{thm}


Алгебра векторных полей на однородной 


поверхности вида $(2)$ 


имеет одно из пяти нижеследующих матричных представлений\/${:}$


\begin{equation}


\frak g_1 =


\begin{pmatrix}


 \frac1{2} \xi_1 & 0 & 0 & p_1 \\


0 & \frac1{2}\xi_1 + i \xi_2 & 0 & p_2\\ 


2i (\overline p_1 + 2 \alpha p_1) & 2 i \overline p_2 & \xi_1 & q \\


0&0&0&0 


\end{pmatrix}


,


\end{equation}


где $ 0 < \alpha \ne 1/2;$


\begin{equation}


\frak g_2 =


\begin{pmatrix}


( \frac1{\alpha} \overline s - 2 s ) p_1 & -\frac1{\alpha} s \overline p_2


 & 0 & p_1 \\


( \frac1{\alpha} \overline s - 2 s ) p_2 


& B_2 & 0 & p_2\\ 


2i (\overline p_1 + 2 \alpha p_1) & 2 i \overline p_2 & 


 \frac 1{\alpha} ( \overline s p_1 + s \overline p_1 ) & q \\


0&0&0&0 


\end{pmatrix}


,


\end{equation}


где $ 0 < \alpha \ne 1/2 $, $ s \in \Bbb R $ или $ s \in i \Bbb R $,


$ B_2 =( \frac1{ 2 \alpha}\overline s - s ) p_1 +


( \frac1{ 2 \alpha} s - \overline s ) \overline p_1 + i \xi_2;$


\begin{equation}


\frak g_3 =


\begin{pmatrix}


 (\frac1{\alpha} \overline s - 2 s) p_1 & 0 & 0 & p_1 \\


0 & 


\frac1{2\alpha} ( \overline s p_1 + s \overline p_1) + i \xi_2


& 0 & p_2\\ 


2i (\overline p_1 + 2 \alpha p_1) & 2 i \overline p_2 & 


 \frac1{\alpha} ( \overline s p_1 + s \overline p_1) & q \\


0&0&0&0 


\end{pmatrix}


,


\end{equation} 


где $ 0 < \alpha \ne 1/2 $, $ s \in \Bbb C; $


\begin{equation}


\frak g_4 =


\begin{pmatrix}


( \frac1{2\alpha} - 2 ) \beta p_1 


 -\frac \beta{ 2 \alpha} \overline p_1 + 


\frac1{\alpha}p_2 


& \frac1{\alpha}p_1


 & 0 & p_1 \\


 - 2 p_1 - \frac1{\alpha} \overline p_1 


& \frac1{\alpha}p_2 & 0 & p_2\\ 


2i (\overline p_1 + 2 \alpha p_1) & 2 i \overline p_2 & 


 \frac 1{\alpha} ( p_2 + \overline p_2 ) & q \\


0&0&0&0 


\end{pmatrix}


,


\end{equation}


где $ \beta \in \Bbb R_+, \ \beta^2 (2 \alpha - 1 ) = 4;$


\begin{equation}


\frak g_5 =


\begin{pmatrix}


A_1 & \frac1{\alpha} ( p_1 + \beta \overline p_2 ) & 0 & p_1 


\\


( - 2 p_1 - \frac1{\alpha} \overline p_1 ) + ( 2 - \frac1{\alpha} ) \beta p_2 


& B_2 & 0 & p_2\\ 


2i (\overline p_1 + 2 \alpha p_1) & 2 i \overline p_2 & C_3 & q \\


0&0&0&0 


\end{pmatrix}


,


\end{equation}


где $\beta \in \Bbb R_+$, $\beta^2 ( 2 \alpha - 1) = 1;$ 


\begin{align*}


A_1 = & 


 - \bigg( 2 + \frac1{\alpha} \bigg) \beta p_1 


 -\frac {2\beta} { \alpha} \overline p_1 + 


\frac1{\alpha}( 2 p_2 + \overline p_2 ),


\\


B_2 =&


\bigg( 2 - \frac2{\alpha} \bigg)


\beta p_1 - \frac1{\alpha} \beta \overline p_1 + 


 \frac1{\alpha} p_2 + \frac2{\alpha} \overline p_2,\\


C_3 =&


\frac 3{ \alpha} ( - \beta ( p_1 + \overline p_1 ) + 


 p_2 + \overline p_2 ).


\end{align*}


\end{thm}





\begin{notenonum}


Во всех алгебрах (9)--(13) имеются 


свободные параметры


$ p_1$, $p_2$, $q $, которые отвечают за сдвиг вдоль поверхности. 


Две последние


алгебры из нашего списка содержат только эти параметры и являются 5-мерными.


Алгебра (9) является 7-мерной и содержит, кроме набора 


$ p,\overline p$, $q $, еще два вещественных параметра $ \xi_1$, $\xi_2 $.


Формулы (10), (11) содержат по 6 вещественных параметров.


\end{notenonum}








При доказательстве теоремы 1 естественно выделить следующий относительно 


простой случай.





\begin{prop}


Если в уравнении $(2)$ однородной поверхности $ M $ 


многочлен $ F_3(z,\ov z) $ тождественно равен нулю, то $ M $ есть квадрика 


$$


 v = ( | z_1 |^2 + |z_2 |^2 ) + \alpha_1 ( z_1^2 + \overline z_1^2 ),


$$


а алгебра векторных полей на ней имеет вид $(9)$.


\end{prop}





 В остальных случаях при изучении равенств (7), (8) 


приходится подробно рассматривать многочлены 


$ F_3 $ и $ F_4 $. Например, пусть ненулевой многочлен $ F_3 $ имеет разложение


$$ 


 F_3 = F_{30} + F_{21} + F_{12} + F_{03},


$$ 


где


\begin{align*}


F_{30} =&


 h_{30} z_1^3 + h_{21} z_1^2 z_2 + h_{12} z_1 z_2^2 + h_{03} z_2^3, \\


F_{21} = &


 (g_{20} z_1^2 + g_{11} z_1 z_2 + g_{02} z_2^2) \overline z_1 +


 (f_{20} z_1^2 + f_{11} z_1 z_2 + f_{02} z_2^2) \overline z_2.


\end{align*}





Из вещественности функции $ F_3 $ следует, что $ F_{12} $ и $ F_{03} $ 


комплексно сопряжены многочленам $ F_{21} $ и $ F_{30} $ соответственно.


Тогда весь многочлен $ F_3 $ определяется десятью коэффициентами:


$h_{30}$, $h_{21}$, $h_{12}$, $h_{03}$, $g_{20}$,


$g_{11}$, $g_{02}$, $f_{20}$, $f_{11}$, $f_{02}$. 





Из тождества веса 2 выводятся равенства


$$


 h_{12} = h_{03} = g_{02} = f_{02} = 0, \ \


 h_{30} = \frac13 ( 4 \alpha g_{20} - \overline g_{20}), 


\ \


 h_{21} = 2 \alpha g_{11} - \overline f_{20}


$$


и доказывается, что все коэффициенты поля $ Z $ на $ M $ выражаются линейно


через $ p$, $\overline p$, $q $ 


и еще два параметра,


аналогичных $ \xi_1$, $\xi_2 $ из (9). Но при


$ F_3(z,\ov z) \ne 0 $ эти параметры не могут быть 


одновременно свободными. В тождестве


(8) веса 3 либо один из них, либо оба выражаются через набор 


$ p$, $\overline p$, $q $.





В дальнейших обсуждениях важную роль играет набор коэффициентов


$ ( h_{21}, g_{11}, f_{20} )$.


При этом оказывается справедливым





\begin{prop}


Для всех однородных поверхностей вида $(2)$ 


выполняется равенство 


$$


 g_{11} = 0 .


$$


\end{prop}





Если коэффициент $ f_{20} $ не 


равен нулю, то растяжением переменных $ z_1$, $z_2 $ его можно превратить в 


единицу. Поэтому изучим случаи $ f_{20} = 0 $ и $ f_{20} = 1 $.





\begin{prop}


Если $ f_{20} = 0 $ для однородной 


поверхности $(2)$, то 


$f_{11} (f_{11} - 2 g_{20} ) = 0$.


При этом равенства 


$ f_{20} = 0$, $f_{11} = 0$ 


дают матричное представление поля $ Z $ в виде $(10)$, а при 


$ f_{20} = 0$, $f_{11} \ne 0 $ 


все поля на $ M $ имеют представление $(11)$.


\end{prop}


 





 Вычисление коммутатора $ [Z_1, Z_2] = Z_1 Z_2 - Z_2 Z_1 $ 


двух матриц вида (10) или (11) показывает, что множества матриц такого вида 


образуют 6-мерные алгебры. 





Завершает доказательство теоремы 1





\begin{prop}


Если $ f_{20} = 1 $ для однородной 


поверхности $(2)$, то 


$$


 f_{11} (f_{11} + 2 g_{20} ) = 0.


$$





 При этом равенства 


$ f_{20} = 1$, $f_{11} = 0 $


дают матричное представление поля 


$ Z $ в виде $(12)$, а при 


$ f_{20} = 1$, $f_{11} \ne 0 $ 


все поля на $ M $ имеют представление $(13)$.


\end{prop}





\begin{notenonum}


 Несложные матричные алгебры можно 


проинтегрировать посредством вычисления матричных экспонент и переходом 


к соответствующим матричным группам Ли. Однородная поверхность является при 


этом орбитой какой-либо точки под действием получаемой группы. 


 Такой подход реализован в работе [10]. Для рутинных вычислений 


в ней использован пакет MAPLE. 





В нашей ситуации вычисление экспонент для относительно просто устроенных


алгебр (10) и (11) дает семейства однородных поверхностей


$$


(1 + v )^t = | 1 + z_1 |^2 \pm | z_2 |^2, \quad t \in \Bbb R \setminus 


\{0,1,2 \}, 


$$


и


$$


v = |z_2|^2 + ( ( 1 + z_1)^A \overline {( 1+ z_1)^A } - 1 ), \quad


A \in \Bbb C, 


$$


соответственно.





Более сложные алгебры (12) и (13) будем 


далее называть алгебрами типов (A) и (B) соответственно. Их интегрирование


основано не на вычислении матричных экспонент, а на решении


систем уравнений в частных производных.


\end{notenonum}





\begin{thm}


Однородная поверхность $ M $, 


имеющая алгебру линейных векторных полей типа {\em (A)}, задается 


$($с точностью до линейных преобразований$)$


уравнением


$$


v = ( |z_1|^2 - |z_2|^2 )^{ \mu}\cdot | z_1 + z_2 |^{ 2 ( 1 - \mu ) }, 


\quad \mu < 0.


$$


\end{thm}








\begin{thm}


Однородная поверхность $ M $, 


имеющая алгебру линейных векторных полей типа {\em (B)}, задается 


$($c точностью до линейных преобразований$)$ 


уравнением одного из двух видов${:}$


$$


(\mathrm B1): v = P_3 ( x,y ) + E |P_2(x,y)|^{3/2}, \quad \alpha \ne 1, 


$$


или


$$


(\mathrm B2): v = \frac{ P_4(x,y)}{1- 2x_1}, \quad \alpha = 1.


$$





Здесь константа $ E $ и 


многочлены $ P_2(x,y)$, $P_3(x,y)$, $P_4(x,y) $ 


от переменных 


$ x_1, y_1$, $x_2$, $y_2 $ однозначно определяются значением параметра 


$ \alpha $. 


\end{thm}





Подробные доказательства теорем 2 и 3, а также точные выражения для 


$ E$, $P_2(x,y)$, $P_3(x,y)$, $P_4(x,y) $ будут выписаны в следующих


двух параграфах статьи. Здесь же сформулируем утверждение, вытекающее 


из них в качестве следствия.





\begin{thm}


Аффинно-однородная вещественная гиперповерхность


$ M $ в $ \Bbb C^3 $, допускающая представление $(2)$, аффинно эквивалентна


одной из поверхностей следующего вида\/${:}$


\begin{gather*}


v = ( | z_1 |^2 + |z_2 |^2 ) + \alpha ( z_1^2 + \overline z_1^2 ),\\


v = |z_2|^2 + ( ( 1 + z_1)^A \overline {( 1+ z_1)^A } - 1 ), \quad


A \in \Bbb C, \\


(1 + v )^t = | 1 + z_1 |^2 \pm | z_2 |^2, \quad t \in \Bbb R \setminus 


\{0,1,2 \}, \\


v = 


( |z_1|^2 - |z_2|^2 )^{\mu} | z_1 + z_2 |^{ 2 (1-\mu)},


\quad \mu < 0,\\


v = 


P_3 ( x,y ) + E |P_2(x,y)|^{3/2}, \quad \alpha \ne 1,\\


v = 


\frac{ P_4(x,y)}{1- 2x_1}, \quad \alpha = 1.


\end{gather*}


\end{thm}








\section{Интегрирование алгебр типа (A)}








Рассмотрим алгебру (A) линейных векторных полей в $ \Bbb C^3 $.


Полагая набор параметров $ ( p_1,p_2,q) $ равным 


$ ( 1,0,0) $,


$ ( i,0,0) $,


$ ( 0,1,0) $,


$ ( 0,i,0) $,


$ ( 0,0,1) $, 


получим пять базисных полей этой алгебры


\begin{align} 


 Z_1 &= 


\left( -2\beta z_1 + \frac1{\alpha} z_2 + 1 \right) 


{\partial}_{z_1} 


- \frac{1 + 2\alpha}{\alpha} 


{\partial}_{z_2} 


+ 2 i (1 + 2\alpha) z_1 {\partial}_{w},\notag  \\


%


 Z_2 &= 


i \left( \frac{ ( 1 - 2\alpha) \beta }{\alpha} 


 z_1 + \frac1{\alpha} z_2 + 1 \right) 


{\partial}_{z_1} 


+ i  \frac{1 - 2\alpha}{\alpha} z_1


{\partial}_{z_2} 


+ 2 (1 - 2\alpha) z_1 {\partial}_{w}, \notag \\


%


 Z_3 &= 


 \frac1{\alpha} z_1 


{\partial}_{z_1} +


 \left( \frac1{\alpha} z_2 + 1 \right) 


{\partial}_{z_2} 


+ 2 ( i z_2 + \frac1{\alpha} w ) {\partial}_{w}, \\


%


 Z_4 &= 


 i \frac1{\alpha} z_1 


{\partial}_{z_1} +


 i \left( \frac1{\alpha} z_2 + 1 \right) 


{\partial}_{z_2} 


+ 2 z_2 {\partial}_{w}, \notag \\


%


 Z_5 &= {\partial}_{w}.\notag


\end{align}


Для каждого из таких полей справедливо тождество (4), т.\,е.


$$


 \re \left( Z_k(\Phi) \right)_{|M} \equiv 0,


$$


где 


$\Phi = \Phi( z,\overline z, u,v ) = - v + F(z,\overline z) $


--- определяющая функция поверхности $ M $.





Ясно, что для такой функции 


$  \Phi_w = i/2$, 


а равенство (4) для поля $ Z_5 $ выполняется автоматически.


 Далее будем действовать только с четырьмя содержательными уравнениями 


(4), соответствующими полям $ Z_1$, $Z_2$, $Z_3$, $Z_4 $ из (14). 





Складывая эти равенства 


попарно первое со вторым, а третье с четвертым


и используя 


очевидные тождества


$$


 a = \re  a - i \re ( i a ), \ \ \


 \overline a = \re  a + i \re ( i a ),


$$


получим 





\begin{prop}


Система четырех уравнений 


$ \re ( Z_k(\Phi))_{|M} = 0 $ для полей 


$Z_k$, $k=\ov{1,4}$, из $(14)$ 


равносильна системе двух комплексных уравнений


\begin{align}


z_1 { F }_{z_1} &+ ( z_2 + \alpha )  F_{z_2} =


F + \alpha \overline z_2, \\ 


z_1 F_{ z_2} &=


\re \left( \frac 1{ 1 + 2\alpha} (- 2 \alpha \beta z_1 + z_2 + \alpha ) 


 F_{z_1} \right ) +\notag\\


&\quad + i


\re \left( \frac i{ 1 - 2\alpha} (( 1 - 2 \alpha) \beta z_1 + z_2 + \alpha )


 F_{z_1} \right ) - \alpha z_1.


\end{align}


\end{prop}





Для упрощения уравнений (15), (16) положим 


$$


H = F + (\alpha z_2 + \alpha \overline z_2 + \alpha^2 )


$$


и 


введем следующие обозначения:


\begin{gather*}


\lambda_1 = \frac {\beta}2, \quad


\lambda_2 = \frac {\alpha \beta}{1 + 2\alpha}, \quad


\lambda_3 = \frac {\beta}{ 2 (1 + 2\alpha)}, \\


l_1 = z_1 - \lambda_1 ( z_2+ \alpha), \ \


l_2 = z_1 - \lambda_2 ( z_2+ \alpha), 


\ \


l_3 = \lambda_3 ( z_2+ \alpha).


\end{gather*}





Отметим здесь, что при любых $ \alpha, \beta $ справедливо равенство


$\lambda_3 = \lambda_1 - \lambda_2$,


а в силу условия-связки 


$ ( 2 \alpha - 1) \beta^2 = 4$, 


означающего, в частности, что


$ \alpha > 1/2 $, 


справедливо и неравенство $ \lambda_2 > \lambda_3 $. 





Учитывая также вещественность функций 


$ F (z, \overline z ) $, $ H (z, \overline z ) $ и вытекающее отсюда 


равенство


$$


 F_{\overline z_1} =


 \overline {F_{z_1} },


$$


перепишем уравнения (15), (16) в виде системы


\begin{equation}


\begin{split}


& z_1 H_{z_1} + ( z_2 + \alpha )  H_{z_2} = H,\\


& z_1 H_{z_2} =( - 


( \lambda_1 + \lambda_2 ) z_1 + \lambda_1 \lambda_2 ( z_2 + \alpha) ) 


 H_{z_1} +


\lambda_3 \overline {(z_1 - \lambda_1 z_2 ) } 


 H_{\overline z_1}.


\end{split}


\end{equation}





Домножим второе из полученных уравнений на 


$ (z_2 + \alpha) $


и результат вычтем из первого уравнения (17), умноженного на $ z_1 $. 


Так придем к системе


\begin{align}


& z_1 H_{z_1} + ( z_2 + \alpha )  H_{z_2} = H,\\


&l_1 l_2 H_{z_1} +l_3 \overline l_1  H_{\overline z_1} 


= z_1 H.\notag


\end{align}


Заметим, что системы (17) и (18) эквивалентны 


при $ \alpha > 0 $ и малых $ z_2 $, т.\,к. при этих условиях сумма


$ (z_2 + \alpha) $ отлична от нуля. Если, кроме того,


и $ z_1 $ близко к нулю, то справедливо неравенство


$$ 


 |l_2|^2 - |l_3|^2 \ne 0. 


$$





\begin{prop}


Для вещественной функции 


$ H(z,\overline z) $


второе уравнение системы $(18)$ равносильно при малых $ z_1$, $z_2 $ 


уравнению


\begin{equation}


 H_{z_1} =


\frac { z_1 \overline l_2 - \overline z_1 l_3 }


{ l_1 ( |l_2|^2 - |l_3|^2 ) } H.


\end{equation}


\end{prop}





Для доказательства этого предложения заметим, что обсуждаемое уравнение 


имеет в краткой форме вид


$ a \zeta + b \overline {\zeta} = c. $


 Отсюда легко получить выражение для 


$ \zeta = H_{z_1} $ 


через $a$, $b$, $c$ и вывести 


(с учетом вещественности $ H(z,\overline z) $) 


формулу, имеющую вид (19). 








Обратимся теперь к системе 


\begin{align}


&z_1 H_{z_1} + ( z_2 + \alpha )  H_{z_2} = H,\\


& H_{z_1} =


\frac { z_1 \overline l_2 - \overline z_1 l_3 }


{ l_1 ( |l_2|^2 - |l_3|^2 ) } H.\notag


\end{align}





\begin{prop}


Любое аналитическое решение первого из 


уравнений системы $(20)$ имеет вид 


\begin{equation}


 H = ( z_2 + \alpha ) g \left(\frac {z_1}{ z_2 + \alpha} \right),


\end{equation}


где $ g(t) $ --- произвольная аналитическая функция одного аргумента.


\end{prop}





Действительно, два первых интеграла системы 


$$


 \frac {d z_1}{z_1} = 


 \frac {d z_2}{z_2 + \alpha} = 


 \frac {d H}H,


$$ 


имеют вид


$$


 \frac { z_1}{z_2 + \alpha} = C_1, \quad 


 \frac H{z_2 + \alpha} = C_2.


$$


Поэтому решение $ H $ исходного уравнения в частных производных 


удовлетворяет равенству


$$


 \frac H{z_2 + \alpha} = g \bigg(\frac { z_1}{z_2 + \alpha} \bigg),


$$


т.\,е. совпадает с (21).





Подставим (21) во второе уравнение сиcтемы (20) и положим


$$ 


 t = \frac { z_1}{z_2 + \alpha}. 


$$


Таким образом, получается уравнение


$$


 g'(t) = 


\frac { ( z_1 \overline l_2 - \overline z_1 l_3 )( z_2 + \alpha ) g(t)}


{ l_1 ( |l_2|^2 - |l_3|^2 ) }


$$


или после деления на


$ ( z_2 + \alpha )^2 \overline{( z_2 + \alpha ) }$


числителя и знаменателя правой части 


\begin{equation}


g'(t) = \frac { A t + B }{ ( t - \lambda_1 )( At - C )} g(t),


\end{equation}


где


$$


 A =  \overline t - \lambda_2, \quad 


 B = - \lambda_3 \overline t, \quad 


 C = \lambda_3^2 + A \lambda_2.


$$





\begin{prop}


Решение уравнения $(22)$ имеет вид


$$ 


 g(t) =


\frac{\lambda_2}{\lambda_1 - \lambda_2} 


\ln( t - \lambda_1 )+


\frac{\lambda_1 - 2 \lambda_2}{\lambda_1 - \lambda_2} 


\ln( At - C ) + D, 


$$


где $ D $ --- произвольная константа.


\end{prop}





Доказательство этой формулы связано с чисто техническими процедурами при 


выражении одних констант через другие в процессе вычисления интеграла от 


рациональной функции.





\begin{notenonum}


С учетом введенных выше обозначений выражение


$$ At - C = At - ( \lambda_3^2 + A \lambda_2 ) = A ( t- \lambda_2) - 


 \lambda_3^2 = 


| t - \lambda_2 |^2 - \lambda_3^2


$$


является вещественным.


\end{notenonum}





Вводя еще одно обозначение 


$$ 


\mu = \frac{\lambda_1 - 2 \lambda_2}{\lambda_1 - \lambda_2} 


$$


и учитывая предложение 7, получаем решение системы (20) в виде


$$


 H(z,\overline z) = 


( z_2 + \alpha ) \bigg( 


\frac{z_1}{z_2 + \alpha} - \lambda_1 


\bigg)^{1-\mu} 


\bigg( \bigg| 


\frac{z_1}{z_2 + \alpha} - \lambda_2 \bigg|^2 


- \lambda_3^2 \bigg)^{\mu} e^D .


$$





Заметим, что систему (20) можно решить лишь с точностью до множителя


$ \varphi (\overline z) = 


\varphi (\overline z_1, \overline z_2)$.


Тогда требование вещественности $ H(z, \overline z) $


означает, что множитель $ e^D $ превращает комплексное выражение


$$


( z_2 + \alpha ) \bigg( 


\frac{z_1}{z_2 + \alpha} - \lambda_1 


\bigg)^{1-\mu} 


$$


в вещественное. Следовательно, 


$$


 e^D = \overline{


( z_2 + \alpha ) \bigg( 


\frac{z_1}{z_2 + \alpha} - \lambda_1 


\bigg)^{1-\mu} } \cdot r,


$$


где $ r $ --- вещественная константа. Тогда 


$$


H(z,\overline z) = r \cdot 


| z_2 + \alpha |^2 \bigg| 


\frac{z_1}{z_2 + \alpha} - \lambda_1 


\bigg|^{2(1-\mu)}


\bigg( \bigg| 


\frac{z_1}{z_2 + \alpha} - \lambda_2 \bigg|^2 - \lambda_3^2 \bigg)^{\mu}


$$


или


$$


H = r \cdot 


 | z_1 - \lambda_1 (z_2 + \alpha) |^{2(1-\mu)}


\left( | z_1


-\lambda_2 ( z_2 + \alpha)|^2 - \lambda_3^2 | z_2 + \alpha|^2 \right)^{\mu}.


$$





С точностью до линейных замен координат получаем в силу этого уравнение 


поверхности


\begin{equation}


v =  | z_1 + z_2 |^{2(1-\mu)}


\left( | z_1|^2 - | z_2 |^2 \right)^{\mu}.


\end{equation}





Для проверки аффинной однородности такой поверхности достаточно заметить, 


что пара эрмитовых форм $ ( H_1, H_2 ) $, где 


$$ 


 H_1 = |z_1|^2 - |z_2|^2 , \quad 


 H_2 = |z_1 + z_2|^2,


$$


сохраняется с точностью до умножения их на некоторые


константы $( c_1,c_2 )$ \ 4-мерной 


группой преобразований 


\begin{equation}


\begin{pmatrix} z_1 \\


 z_2 \end{pmatrix} \to


\lambda e^{i\theta} 


\begin{pmatrix} \sqrt{ 1+t^2} + is & t + is \\


 t - is & \sqrt{ 1+t^2} - i s 


\end{pmatrix}


\begin{pmatrix} z_1 \\


 z_2 \end{pmatrix}, 


\end{equation}


где $ (s,t,\lambda,\theta) \in \Bbb R^4$.


 При этом действие группы (24) в пространстве $ \Bbb C^2 $ является 


транзитивным вблизи точки 


$ (z_1, z_2) = ( - \lambda_2 \alpha, - \lambda_3 \alpha ) $.





 Дополняя эти преобразования сдвигами в направлении оси $ u $, 


получаем 5-мерную группу аффинных преобразований пространства $ \Bbb C^3 $, 


транзитивно действующую на многообразии (23).





\begin{notenonum}


Из тех же соображений следует аффинная однородность 


поверхностей, задаваемых уравнениями 


$$


v = 


 | z_1 + z_2 |^{\nu}


\left( | z_1|^2 - | z_2 |^2 \right)^{\mu}


$$


с произвольными вещественными показателями $ \mu $ и $ \nu $. 


\end{notenonum}








\section{Интегрирование алгебр типа (B)}








Базисные поля $ Z_1$, $Z_2$, $Z_3$, $Z_4$, $Z_5 $ алгебры (13), 


соответствующие наборам $ (p_1,p_2,q) $ вида 


$ (1,0,0) $, 


$ (i,0,0) $, 


$ (0,1,0) $, 


$ (0,i,0) $, 


$ (0,0,1) $, 


имеют достаточно сложный вид. 


Перейдем от этих полей к новому базису


$$


 Z_1 + \beta Z_2, \quad \beta Z_2 + Z_4, \quad \alpha \beta Z_3, 


\quad \beta Z_4, \quad Z_5 


$$


изучаемой алгебры, сохраняя для него старые обозначения. 





 С учетом формулы-связки 


$ ( 2 \alpha - 1 ) \beta^2 = 1$ систему уравнений 


$$


 \re (Z_k(\Phi))_{|M} = 0, \quad k=\ov{1,4},


$$


можно тогда записать в виде 


\begin{gather}


 \re \left( ( -2 \beta z_1 + 2 \beta^2 z_2 + 1) 


 F_{z_1} + ( -2 z_1 + 2 \beta z_2 + \beta) 


 F_{z_2} \right) -\notag \\


- \re ( ( 1 + 2 \alpha ) z_1 + \beta z_2 ) = 0,\notag \\


\re \left( 


 i \beta z_1  F_{z_1} +  i F_{z_2} \right) +


 \re \left(i 


( - \tfrac 1{\beta} z_1 + z_2 ) \right) = 0, \\


 \re \left( 


( 3 \beta z_1 + \beta^2 z_2 )  F_{z_1} + 


( z_1 + 3 \beta z_2 + \alpha \beta)  F_{z_2} \right) + \notag \\


+ \re ( - \alpha \beta z_2 + 3 i \beta (u + i F ) ) = 0,\notag \\


 \re \left( 


 i \beta ( z_1 - \beta z_2 )  F_{z_1} + 


( i z_1 - i \beta z_2 + i \alpha \beta)  F_{z_2} \right) +


 \re ( i \alpha \beta z_2 ) = 0.\notag


\end{gather}





На следующем этапе сделаем замену переменных.





\begin{prop}


При замене переменных


$$ 


 \zeta_1 = z_1 - \beta z_2, \ \ 


 \zeta_2 = z_1 + \beta z_2


$$


система уравнений $(25)$ примет вид


\begin{gather}


 \re \left( ( 1- \beta^2) 


F_{\zeta_1} + ( 2 \alpha \beta^2 - 4 \beta \zeta_1 ) 


F_{\zeta_2} \right)= 


\alpha \re \zeta_1 + (1 + \alpha) \re \zeta_2, \notag \\


 \re \left( 2 i \beta 


F_{\zeta_2 } \right)= - \frac 1{\beta} \im \zeta_1, \\


 \re \left( \beta ( 2 \zeta_1- \alpha \beta ) 


F_{\zeta_1} + \beta ( 4 \zeta_2 + \alpha \beta ) F_{\zeta_2} 


\right)= 3 \beta F + \frac{\alpha}2 \re (- \zeta_1 + \zeta_2 ),\notag\\


 \re \left( - i \alpha \beta^2


F_{\zeta_1} +  i \beta ( 2 \beta \zeta_1 + \alpha \beta^2 ) 


F_{\zeta_2} \right)= 


\frac{\alpha}2 \im (- \zeta_1 + \zeta_2 ).\notag


\end{gather}


\end{prop}





Для дальнейшего упрощения системы (26) заметим, что здесь, в отличие от


интегрирования алгебр типа (A), удобно перейти к вещественным координатам.


Полагая 


$$


\xi_k = \re \zeta_k, \quad \eta_k = \im \zeta_k, \quad k = 1,2,


$$


получаем в качестве следствия из предложения 9 следующую систему уравнений 


в частных производных:


\begin{gather}


 ( 1- \beta^2) 


F_{\xi_1} + ( 2 \alpha \beta^2 - 4 \beta \xi_1 ) 


F_{\xi_2} - 4 \beta \eta_1F_{\eta_2} 


= 2 \alpha \xi_1 + 2 ( 1 + \alpha) \xi_2, \notag \\


F_{\eta_2 } = - \frac 1{\beta^2} \eta_1, \\


 ( 2 \beta \xi_1- \alpha \beta^2 ) F_{\xi_1} + 


 2 \beta \eta_1 F_{\eta_1} 


 + ( 4 \beta \xi_2 + \alpha \beta^2 ) F_{\xi_2} +\notag \\


 + 4 \beta \eta_2 F_{\eta_2} = 


6 \beta F + \alpha (- \xi_1 + \xi_2 ), \notag \\


 \alpha \beta^2F_{\eta_1} +  2 \beta \eta_1


F_{\xi_2}  - ( 2 \beta \xi_1 + \alpha \beta^2 ) 


F_{\eta_2} = 


\alpha (\eta_1 - \eta_2 ).\notag


\end{gather}





Самое простое уравнение этой системы имеет решение


$$


 F (\xi_1, \eta_1, \xi_2, \eta_2 ) = 


- \frac 1{\beta^2} \eta_1 \eta_2 + H( \xi_1, \xi_2, \eta_1 ), 


$$


где $ H(\xi_1, \xi_2, \eta_1 ) $ --- произвольная аналитическая функция 


от трех вещественных переменных. 





Подставляя эту формулу в три оставшиеся уравнения системы (27), получим ее


сокращенный вариант


\begin{gather}


( 1- \beta^2) 


H_{\xi_1} + ( 2 \alpha \beta^2 - 4 \beta \xi_1 ) 


H_{\xi_2} = 


2 \alpha \xi_1 + 2 ( 1 + \alpha) \xi_2 - \frac4{\beta} \eta_1^2,\notag\\


 ( 2 \beta \xi_1- \alpha \beta^2 ) 


H_{\xi_1} +  2 \beta \eta_1 H_{\eta_1} +


( 4 \beta \xi_2 + \alpha \beta^2 ) H_{\xi_2} = 


6 \beta H + \alpha (- \xi_1 + \xi_2 ), \\


 \alpha \beta^2H_{\eta_1} +  2 \beta \eta_1H_{\xi_2} = 


- \frac 2{\beta} \xi_1 \eta_1, \notag


\end{gather}


содержащий три уравнения относительно производных функции $ H $ 


по переменным $ \xi_1$, $\xi_2$, $\eta_1 $. 


 


Для дальнейшего упрощения полученной системы удобно еще раз произвести


замену переменных.





\begin{prop}


При замене переменных


$$ 


t_1 = \xi_1, \quad t_2 = \eta_1^2 - \alpha \beta \xi_2, \quad t_3 = \eta_1


$$


система уравнений $(28)$ примет вид


\begin{gather}


( 1- \beta^2) H_{t_1} - 


2 \alpha \beta^2 ( \alpha \beta - 2 t_1 ) 


H_{t_2} = 


2 \alpha t_1 - \frac 2{\alpha \beta} ( 1 + \alpha) t_2 


+ 2\frac {1 - \alpha}{\alpha \beta} t_3^2,\notag \\


 ( 2 \beta t_1 - \alpha \beta^2 ) H_{t_1} 


 + ( 4 \beta t_2 - \alpha^2 \beta^3 ) 


H_{t_2}  + 2 \beta t_3 H_{t_3} = 


6 \beta H - \alpha t_1 + \frac 1{ \beta} (t_3 - t_2 ), \notag \\


\alpha \beta^2 H_{t_3} =


- \frac 2{\beta} t_1 t_3.


\end{gather}


\end{prop}





Решением последнего уравнения этой системы является


$$


 H = - \frac{\alpha}{\beta } t_1 t_3^2 + g( t_1, t_2 ),


$$ 


где $ g(t_1,t_2) $ --- произвольная аналитическая функция двух переменных. 





 Подставляя эту формулу в остальные уравнения системы (29), получим


систему двух уравнений относительно производных функции $ g $ 


по переменным $ t_1$, $t_2 $


\begin{gather}


 ( 1- \beta^2) g_{t_1} +


 2 \alpha \beta^2 ( 2 t_1 - \alpha \beta )g_{t_2} = 


 2 \alpha t_1 - \frac 2{\alpha \beta} ( 1 + \alpha) t_2, \\


%


 ( 2 \beta t_1 - \alpha \beta^2 ) g_{t_1} 


 + ( 4 \beta t_2 - \alpha^2 \beta^3 ) g_{t_2} = 


6 \beta g - ( \alpha t_1 + \frac{t_2}{\beta} ) .


\end{gather}








В зависимости от 


значения параметра $ \beta $ (и связанного с ним $ \alpha $) при рассмотрении


системы (30), (31) естественно выделить два случая.








1. В случае $ \beta^2 = 1$ $(\alpha = 1) $


существенно упрощается уравнение (30). Его решение 


имеет вид


$$


 g(t_1,t_2) = \frac{ t_1 t_2 - t_2^2 }{ 2 t_1 - 1 } + 


\varphi(t_1),


$$


где $ \varphi(t_1) $ --- произвольная аналитическая функция одного аргумента.


 Тогда из (31) получаем обыкновенное дифференциальное уравнение


$$


 (2 t_1 - 1) \varphi'(t_1) - 6 \varphi(t_1) = 


 \frac{ 2 ( t_1 - t_1^2 ) }{ ( 2 t_1 - 1 )^2 }


$$


относительно функции $ \varphi(t_1) $.


 Решая это уравнение и возвращаясь к переменным 


$ \xi, \eta $, получаем уравнение искомой однородной поверхности в виде


$$


 v = \frac{ 


 (\xi_1^2 + \eta_1^2 )^2 - 2 \xi_1 ( \xi_1^2 + \eta_1^2 ) -


 2 \xi_2 \eta_1^2 + \xi_1^2 + \xi_1 \xi_2 + \xi_2^2}{1-2 \xi_1} . 


$$


 


2. В случае $ \beta^2 \ne 1$ $(\alpha \ne 1) $ 


начнем изучение системы (30), (31) с ее второго уравнения.





\begin{prop}


Решение уравнения $(31)$ имеет вид


\begin{equation}


 g( t_1,t_2) = 


\frac 1{\alpha^2} + \frac{\alpha}{4 \beta} ( t_1 - A ) +


\frac 1{2 \beta^2} ( t_2 - A^2 ) +


 ( t_1 - A )^3 \varphi \left


( \frac {t_2 - A^2}{ ( t_1 - A)^2 } \right),


\end{equation}


где 


$ A = \alpha \beta/2$, $\varphi $ --- произвольная аналитическая 


функция одного вещественного переменного. 


\end{prop}





Для доказательства разделим уравнение (31) на $ 2 \beta $ и запишем в виде 


$$


 ( t_1 - A ) \frac{\partial g}{\partial t_1} 


 + 2 ( t_2 - A^2 ) \frac{\partial g}{\partial t_2} 


= 3 g - ( B t_1 + D t_2 ), 


$$


где $ B = \alpha / ( 2 \beta)$, $D = 1 / ( 2 \beta^2)$.


Соответствующая этому уравнению система


обыкновенных дифференциальных уравнений имеет вид


$$


 \frac{ d t_1}{ t_1 - A} =


 \frac{ d t_2}{ 2( t_2 - A^2 )} =


 \frac{ d g}{ 3 g - ( B t_1 + D t_2 ) }.


$$


Два независимых первых интеграла этой системы 


$$


C_1 = \frac {t_2 - A^2}{ ( t_1 - A)^2 },\quad 


C_2 = \frac {g - N/3 - (B/2)( t_1 - A) - D( t_1 - A)^2 } { (t_1 - A)^3 },


$$


где $ N = D A^2 + B A = 3 \alpha^2/8$,


приводят к формуле (32).





После подстановки (32) в уравнение (30) получим для функции $\varphi(s)$


обыкновенное дифференциальное уравнение


\begin{equation}


( s - \lambda)\varphi'(s) = \frac 32 \varphi + \mu s,


\end{equation}


где 


$$


\lambda = \frac{ 1 + \beta^2 }{ 1 - \beta^2 }, \quad


\mu = \frac{ 1 + 3 \beta^2 }{ \beta ( 1 - \beta^4 ) }.


$$





Заметим, что введенный параметр $ \lambda $ может принимать 


значения (в зависимости от 


$ \alpha $ или, что то же самое, от $ \beta $) 


из объединения интервалов


$ ( - \infty, -1 ) \cup (1,\infty )$.


 Переменная же 


$$ 


 s = \frac {t_2 - A^2}{ ( t_1 - A)^2 } 


$$ 


близка при малых $ t_1$, $t_2 $ к $ - 1 $.


Поэтому, интегрируя уравнение (33), получаем 


\begin{equation}


 \varphi(s) = E | s - \lambda |^{3/2} - 2 \mu s + \frac 43 \lambda \mu,


\end{equation}


где модуль раскрывается по-разному в зависимости от значения 


$ \alpha $.


 


Подставляя (34) в (32) и учитывая предыдущие 


интегрирования, получаем окончательно уравнение


\begin{multline*}


v = 


 \frac {4( 1+ 3 \beta^2) }{ 3\beta( 1- \beta^2)^2} \xi_1^3 - 


 \frac 4{ \beta( 1- \beta^2)} \xi_1 \eta^2_1 


- \frac 1{\beta^2} \eta_1 \eta_2 + 


 \frac { 2 \alpha }{ 1 - \beta^2} \eta_1^2 -  \\


%


- \frac { 2 \alpha( 1+ 3 \beta^2) }{ ( 1- \beta^2)^2} \xi_1^2 +


 \frac { 2 \alpha - 1 }{ 1- \beta^2} \xi_1 \xi_2 +


\frac { \alpha }{ \beta } \lambda (\lambda \xi_1 - \xi_2 ) -


 \frac { \alpha }{ 6 \beta^2 } \lambda^2 +


\\


+ \frac { 16 }{ 3 \beta ( 1- \beta^2)^2 | \lambda + 1 |^{3/2} } 


 \left| \eta_1^2 - \lambda \xi_1^2 + 2 A ( \lambda \xi_1 - \xi_2) -


 ( \lambda + 1) A^2 \right|^{3/2}


\end{multline*}


однородной поверхности $M$.





Напомним, что все коэффициенты в этой формуле, т.\,е. 


$ \alpha$, $\beta$, $\lambda$, $A $ однозначно выражаются через один параметр


$ \alpha $ или $ \beta $.








\begin{thebibliography}{99}





\bibitem{1}


Широков П.А., Широков А.П. {\em Аффинная дифференциальная геометрия}.


-- М.: Физ\-мат\-лит, 1959. -- 320 c.


\bibitem{2}


Doubrov B., Komrakov B., Rabinovich M. {\em Homogeneous surfaces 


in the $3$-dimensional affine geometry} // Prepr. Ser. Pure


Math. / Inst. Math.


Univ. -- Oslo, 1995. -- \No\,4. -- P.\,1--26. 


\bibitem{3}


Cartan E. {\em Sur la geometrie pseudoconforme des hypersurfaces de deux 


variables complexes} // Ann. Math. Pura Appl. -- 1932. -- V.\,11.


-- \No\,4. -- P.\,17--90


(Oeuvres II, 2, 1231--1304).


\bibitem{4}


Лобода А.В. {\em О размерности группы, транзитивно действующей 


на гиперповерхности в $ \Bbb C^3 $} // Функц. анализ и его прилож. 


-- 1999. -- Т.\,33. -- Вып.\,1. -- C.\,68--71.


\bibitem{5}


Лобода А.В. {\em Однородные строго псевдо-выпуклые гиперповерхности в 


 $ \Bbb C^3 $ с двумерными группами изотропии} // Матем. сб.


-- 2001. -- Т.\,192. -- \No\,12. -- С.\,3--24. 


\bibitem{6}


Chern S.S., Moser J.K. {\em Real hypersurfaces in complex manifolds} //


Acta Math. -- 1974. -- V.\,133. -- \No\,3. -- P.\,219--271.


\bibitem{7}


Белошапка В.К. {\em Однородные гиперповерхности в $ \Bbb C^2 $} 


// Матем. заметки. -- 1996. -- Т.\,60. -- \No\,5. -- С.\,760--764. 


\bibitem{8}


Stanton N.K. {\em Infinitesimal CR automorphisms of rigid hypersurfaces} 


// Amer. J. Math. -- 1995. -- V.\,117. -- \No\,1. -- P.\,141--167. 


\bibitem{9}


Гузеев Р.Н., Лобода А.В. {\em О нормальных уравнениях аффинно-однородных 


выпуклых поверхностей пространства $ \Bbb R^3 $} // Изв. вузов.


Математика. -- 2001. -- \No\,3. -- С.\,25--32. 


\bibitem{10}


Eastwood M., Ezhov V.V. {\em On affine normal forms and a classification


of homogeneous surfaces in affine three-space} // Geom. dedic. -- 1999.


-- V.\,77. -- P.\,11--69.


\bibitem{11}


Лобода А.В., Бугаева Ж.А., Ходарев А.С. {\em О линейной однородности


жестких вещественных гиперповерхностей $3$-мерного комплексного 


пространства} // Тез. докл. школы-семинара по геометрии и анализу, 


посвящ. 90-летию Ефимова Н.В. -- Абрау-Дюрсо, 2000. -- С.\,131--133.





\end{thebibliography}





\vskip 2cm





\postup{\it Воронежский государственный}{\qquad \it Поступили}


\postup{\it архитектурно-строительный}{\it первый вариант $28.06.2001$}


\postup{\it университет}{\it окончательный вариант $05.02.2003$}





\label{end}


\end{document}





